\section{Introduction}
Reaction--diffusion models describe an expanding population through spatial dispersal and local growth. The Fisher--KPP equation emerged from Fisher and from Kolmogorov, Petrovskii and Piskunov \citep{fisher1937,kpp1937}. Related proliferation--invasion formulations describe glioma-cell migration through anatomically heterogeneous brain tissue. Distinguishing white- and gray-matter motility permits predictions of preferential spread even when local transport is isotropic \citep{swanson2000}.

A single viable-cell field omits microenvironmental processes that can distinguish an infiltrating margin from a poorly oxygenated interior. The histological illustration motivating this work, from Comba et al., emphasizes invasion routes, vascular niches, reactive cells and necrosis \citep{comba2021}; the present model represents only a small subset of that conceptual picture. Earlier mechanistic modeling connected oxygen, cellular phenotypes and necrosis to pseudopalisading patterns \citep{martinez2012}. Our reduction does not reproduce that model or explicitly resolve pseudopalisades. Oxygen-dependent proliferation and transfer into a necrotic compartment introduce additional assumptions and parameters. A more elaborate model can reproduce more patterns without providing uniquely identifiable biological explanations. Model complexity, observability and prediction therefore need to be assessed together.

Recent work illustrates the opportunities and constraints. Hiremath et al. investigated identifiability and selection among glioma treatment models using noisy virtual cohorts \citep{hiremath2025}. Miniere et al. incorporated serial MRI into a two-species model for 21 patients receiving chemoradiation, comparing sequential updates with prediction from limited imaging visits \citep{miniere2025}. Zhang et al. explored physics-informed inference from a single structural MRI snapshot using characteristic nondimensional parameters \citep{zhang2025}. These approaches motivate explicit consideration of what observations constrain; they do not establish that every mechanistic rate is identifiable from one scan.

Measurement validity is an additional issue. Independent histological evaluation of MRI-derived cellularity maps in 2025 found different performance in treatment-naive and recurrent gliomas \citep{nocera2025}. A 2026 test--retest study combined diffusion, perfusion and oxygen-enhanced MRI to construct repeatable habitats in a murine glioma model \citep{das2026}. Repeatability supports longitudinal inference, while biological specificity requires separate evidence. Recent 2026 preprints investigate poromechanical forecasting, sequential biomechanical model selection, and biopsy-derived kinetic information \citep{andaloussi2026,noman2026,hofmann2026}. Their preprint status and experimental scope are retained here.

The present work develops a compact hierarchy from Fisher--KPP to oxygen-limited growth and coupled viable-tumor--oxygen--necrosis dynamics. An additional MRI pilot compares isotropic and dMRI-directed growth on sub-PAT16 anatomy, with paired sensitivity ensembles and measured resting-state connectivity overlays. The workflow provides reproducible anatomical simulations, interactive exploration, numerical assessment, and staged synthetic recovery of homogeneous rates, separate tissue rates and joint tissue migration/anisotropy on PAT16 anatomy, with uncertainty, withheld-scan tests and anatomy sensitivity assessment. We ask how oxygen and necrosis alter the simulated interior relative to threshold-defined extent, which assumptions dominate that change, whether baseline, separate white-/gray-matter rates and directional migration anisotropy can be recovered under controlled observation assumptions, how discretized anatomy affects those estimates, and what measurements would be needed to interpret the parameters biologically. The study does not claim a new biological law or clinical forecasting performance.
